% Source mirror: `paper/tex/design_rationale.tex`.
\section{设计依据}
\label{sec:design-rationale}

\SPHINCSA\ 的设计思路是：沿用原始 \SPHINCSP\ 的总体结构，引入优化后的 \WOTSCS\ 一次性签名方案，将底层少次签名替换为 BPORS+FP，并在超树上采用 Diverse Layer Heights（分层高度多样化）优化。\SPHINCSP\ 的基本出发点如下：可以通过构造一棵巨大的二叉认证树，并为每个待签消息随机选择一个叶子，来建立无状态的哈希签名方案。然而，这样的树必须极高，其高度通常需要达到大约两倍的安全级别，从而导致签名尺寸极不实用。使用超树代替单棵二叉认证树，可以在速度与签名大小之间进行折中，但仅靠这一点仍不足以同时获得实用的尺寸和速度。

\paragraph{数学困难问题。}
\SPHINCSA\ 是一种无状态、抗量子的数字签名方案，其安全性沿用 \SPHINCSP\ 的证明框架，继承自底层哈希与 XOF 原语在面对通用攻击时的困难性。从构造层面看，Kylin 可被视为对 NIST 标准化 SPHINCS+ 设计的改进：它保留基于哈希的超树，将 FORS 替换为 BPORS+FP，同时将 Winternitz 一次性签名层替换为 \WOTSCS。

\paragraph{参数选择准则。}
主参数选择规则是在尽可能接近原始 \SPHINCSP\ 风格实现与分析方式的前提下，优于可比较的 \SPHINCSP\ 风格参数配置。在固定的 NGCC 安全级别要求下，搜索优先级依次为：更小的签名、更快的验证、再到更快的签名，同时保持单密钥签名预算为 $q=2^{64}$。因此，\cref{tab:ngcc-parameters} 中选定的实例主要依赖三个设计杠杆：常和 \WOTSCS\ 编码、BPORS+FP 参数 $(K,B,t,k,m_{\max})$，以及超树的分层高度配置 $(h_0,\ldots,h_{d-1})$。

由于一次性签名在 \SPHINCSP\ 超树结构中被大量使用，优化一次性签名的尺寸会直接影响整体性能。借助 \WOTSCS，Kylin 可以在保持运行时间大致相当的同时，将整个签名方案重新调参为更小的签名。BPORS+FP 进一步通过 bucket thinning 和固定全局认证节点预算减少底层少次签名载荷，并沿用了~\cite{EPRINT:AbrKat25} 中 forced-pruning 式尺寸控制视角。Diverse Layer Heights 进一步消除了 $h$ 必须能被 $d$ 整除这一不必要限制：只要各层高度之和为 $h$，各层就可以不同，因此搜索过程能够使用像 \texttt{Kylin-256} 的 $(5,6,6,6,6,6,6,6,5,6)$ 这类配置。

\paragraph{失败概率与安全影响。}
Kylin 不存在类似解密或封装场景中的正确性失败事件。对于诚实生成的密钥与诚实生成的签名，验证过程是确定性的；除实现故障或底层密码学假设失效外，验证都会接受。因此，参数搜索中使用的概率项是伪造成功概率上界，而不是正确性失败概率。

对于同一密钥对下签名预算为 $q$ 的场景，参数搜索中使用的具体安全指数来自 \cref{eq:proof-generic-prob} 中的 BPORS+FP cover 项。选定的主实例均取 $q=2^{64}$，其安全指数如 \cref{tab:ngcc-parameters} 所示。这些指数用于评估重复底层 key 选择及通用伪造事件对所声明安全级别的影响。

\paragraph{其他考虑。}
仓库中还保留了 $q=2^{24}$ 的附加参考参数配置。这些配置不是 NGCC 主提交通道中的实例，但对于单密钥签名预算较小、同时对签名尺寸更敏感的应用场景仍然有参考价值；参见 \cref{apx:q24-profiles}。

我们能够证明，\WOTSCS\ 中使用的编码方案在所有树结构编码方案中都是尺寸最优的，这个结论甚至覆盖了那些不一定能够高效计算的编码方案。进一步地，就我们目前所知，现有任何基于 DAG 的方案都无法在签名尺寸上优于该方案，因此我们的方案可视为已知方案中的尺寸最优者~\cite{optimalwots}。我们也注意到基于拒绝采样的概率型编码方案~\cite{cryptoeprint:2022/778}，并认为这些技术具有独立研究价值。
